\documentclass[a4paper, serif, 10pt]{moderncv}

%% moderncv
% color and style
\moderncvcolor{black}
\moderncvstyle{banking}

%% page layout/margins
\usepackage[top=2.5cm, bottom=2.5cm, left=1.5cm, right=1.5cm]{geometry}

\nopagenumbers{}

%% Babel config for Dutch language
% ref:
% - https://latex3.github.io/babel/guides/locale-dutch.html
\usepackage[dutch]{babel}

%% fontspec
\usepackage{fontspec}

\IfFontExistsTF{Linux Libertine}{
  \setmainfont[Ligatures={TeX, Common}, Numbers=OldStyle]{Linux Libertine}
}{}
\IfFontExistsTF{Linux Libertine O}{
  \setmainfont[Ligatures={TeX, Common}, Numbers=OldStyle]{Linux Libertine O}
}{}

%% CTeX
% CJK support, used to show CJK characters in the resume
%
% - heading=false: disable ctex heading
% - punct=kaiming: use kaiming punctuations styles for CJK
% - scheme=plain: use plain scheme, do not override \normalsize font size
% - space=auto: space settings for CJK characters
%
% CJK fonts are resolved per locale language, e.g. Noto Serif/Sans CJK JP
% for Japanese resumes, with fallback to the next region if not installed
%
% ref:
% - http://ctan.mirrorcatalogs.com/language/chinese/ctex/ctex.pdf
\usepackage[UTF8, heading=false, punct=kaiming, scheme=plain, space=auto]{ctex}

\IfFontExistsTF{Noto Serif CJK SC}{
  \setCJKmainfont{Noto Serif CJK SC}
}{}
\IfFontExistsTF{Noto Sans CJK SC}{
  \setCJKsansfont{Noto Sans CJK SC}
}{}

%% line spacing
\usepackage{setspace}
\setstretch{1.125}

%% URL styling - use normal text instead of monospace
\urlstyle{same}

%% Patch \httplink and \httpslink to handle full URLs
% The original moderncv macros always prepend http:// or https:// to URLs.
% This causes issues when the URL already has a protocol (e.g., https://example.com)
% resulting in malformed URLs like https://https://example.com.
% This patch checks if the URL already contains "://" and uses it directly if so.
% Note: We use \str_set:Nx (x-expansion) to expand macros like \@homepage first.
\ExplSyntaxOn
\str_new:N \g_yamlresume_url_str

\RenewDocumentCommand{\httpslink}{O{}m}{
  \str_set:Nx \g_yamlresume_url_str {#2}
  \str_if_in:NnTF \g_yamlresume_url_str {://}
    { \url{#2} }
    { \url{https://#2} }
}

\RenewDocumentCommand{\httplink}{O{}m}{
  \str_set:Nx \g_yamlresume_url_str {#2}
  \str_if_in:NnTF \g_yamlresume_url_str {://}
    { \url{#2} }
    { \url{http://#2} }
}
\ExplSyntaxOff

\name{Sanne de Vries}{}
\title{Geregistreerde verpleegkundige met focus op acute en chronische zorg}
\phone[mobile]{+31 6 1234 5678}
\email{sanne.devries@zorgmail.nl}
\homepage{https://www.sanne-de-vries.nl}

\address{Maliebaan 45, Utrecht, Utrecht, Nederland, 3512 JE}{}{}

\extrainfo{{\small \faLinkedin}\ \href{https://www.linkedin.com/in/sanne-de-vries-rn}{@sanne-de-vries-rn} {} {} {} • {} {} {} 
{\small \faMedium}\ \href{https://medium.com/@sannezorg}{@sannezorg}}

\begin{document}

\maketitle

\section{Basis Informatie}

\cvline{}{Geregistreerde verpleegkundige met meer dan zes jaar ervaring in de ziekenhuiszorg en de eerstelijnszorg. Sterk in het leveren van persoonsgerichte zorg, het opstellen en evalueren van zorgplannen en het veilig uitvoeren van medicatie- en wondzorg.
%
\begin{itemize}
\item Ervaren in acute zorg, chronische zorg en palliatieve zorg
\item Vaardig met elektronische patiëntendossiers (HiX, EPIC) en zorgtechnologie
\item Communicatief sterk, sociaal en gericht op samenwerking met patiënten, naasten en multidisciplinaire teams
\item Proactief in kwaliteitsverbetering, patiëntveiligheid en professionele ontwikkeling
\end{itemize}}
\section{Opleidingen}

\cventry{sep 2014–jul 2018}
        {Bachelor, HBO-Verpleegkunde, Score: 8.4}
        {Hogeschool Utrecht}
        {\url{https://www.hu.nl}}
        {}
        {\begin{itemize}
\item Afgestudeerd met een brede basis in ziekenhuiszorg, chronische zorg en preventie
\item Tijdens stages ervaring opgedaan op een chirurgische afdeling en in de wijkverpleging
\item Afstudeeronderzoek gedaan naar medicatieveiligheid bij ouderen
\end{itemize}
\textbf{Cursussen}: Verpleegkundig handelen, Anatomie en fysiologie, Pathologie, Farmacologie, Verpleegkundige ethiek, Communicatie en gespreksvoering, Kwaliteit en veiligheid in de zorg, Zorgtechnologie}
\section{Werkervaring}

\cventry{mrt 2021–Heden}
        {Verpleegkundige}
        {Sint-Maartens Kliniek}
        {\url{https://www.sintmaartenskliniek.nl}}
        {}
        {\begin{itemize}
\item Verantwoordelijk voor de verpleegkundige zorg aan patiënten op de afdeling orthopedie en revalidatie
\item Stelt zorgplannen op, voert deze uit en evalueert ze samen met patiënt en naasten
\item Voert medicatiebeheer, wondzorg en pijnmanagement uit volgens protocollen
\item Werkt samen met artsen, fysiotherapeuten, apothekers en verzorgenden in een multidisciplinair team
\item Draagt bij aan kwaliteitsverbetering door deelname aan audits en het melden van incidenten
\end{itemize}
\textbf{Trefwoorden}: Klinische zorg, Patiëntveiligheid, Zorgplanning, Medicatiebeheer, Wondzorg, Revalidatie}

\cventry{aug 2018–feb 2021}
        {Verpleegkundige}
        {Zorgcentrum Parkzicht}
        {\url{https://www.zorgcentrumparkzicht.nl}}
        {}
        {\begin{itemize}
\item Verleende persoonsgerichte zorg aan ouderen met complexe somatische en psychogeriatrische zorgvragen
\item Coördineerde zorgprocessen, bewaakte veiligheid en ondersteunde verzorgenden bij de uitvoering van zorg
\item Voerde verpleegtechnische handelingen uit, waaronder wondzorg, insulinetoediening en katheterzorg
\item Had aandachtsgebied palliatieve zorg en begeleidde families in de laatste levensfase
\item Nam deel aan multidisciplinair overleg en familiegesprekken
\end{itemize}
\textbf{Trefwoorden}: Ouderenzorg, Palliatieve zorg, Zorgcoördinatie, Psychogeriatrie, Familiebegeleiding}
\section{Talen}

\cvline{Nederlands}{Moedertaal of tweetalige vaardigheid \hfill \textbf{Trefwoorden}: Moedertaal}
\cvline{Engels}{Volledige professionele vaardigheid \hfill \textbf{Trefwoorden}: TOEFL iBT 105, Professionele zorgcommunicatie}
\cvline{Duits}{Beperkte werkvaardigheid \hfill \textbf{Trefwoorden}: Goethe-Zertifikat B1}
\section{Vaardigheden}

\cvline{Klinische Verpleegkunde}{Expert \hfill \textbf{Trefwoorden}: Anamnese, Zorgdiagnose, Zorgplanning, Interventies, Evaluatie}
\cvline{Medicatiebeheer}{Geavanceerd \hfill \textbf{Trefwoorden}: Medicatiecontrole, Toediening, Bijwerkingen, Polyfarmacie, Veiligheid}
\cvline{Wondzorg}{Geavanceerd \hfill \textbf{Trefwoorden}: Wondbeoordeling, Verbandleer, Infectiepreventie, Compressietherapie}
\cvline{Communicatie}{Expert \hfill \textbf{Trefwoorden}: Gespreksvoering, Familiegesprekken, Multidisciplinair overleg, Voorlichting}
\cvline{Digitale Zorg}{Intermediair \hfill \textbf{Trefwoorden}: HiX, EPIC, Nedap, Zorgtechnologie, Microsoft Teams}
\section{Onderscheidingen}

\cventry{nov 2023}
        {Sint-Maartens Kliniek}
        {Verpleegkundige van het Jaar 2023}
        {}
        {}
        {Erkend voor uitmuntende patiëntgerichte zorg, sterke samenwerking en initiatief in het verbeteren van medicatieveiligheid op de afdeling.}
\section{Certificaten}

\cventry{aug 2018}
        {CIBG}
        {BIG-registratie Verpleegkundige}
        {\url{https://www.bigregister.nl}}
        {}
        {}

\cventry{mrt 2024}
        {Nederlandse Hartstichting}
        {Basale reanimatie (BLS)}
        {\url{https://www.hartstichting.nl}}
        {}
        {}

\cventry{apr 2024}
        {Nederlandse Hartstichting}
        {Gevorderde reanimatie (ALS)}
        {\url{https://www.hartstichting.nl}}
        {}
        {}
\section{Publicaties}

\cventry{jun 2023}
        {Medicatieveiligheid bij ouderen in de thuissituatie}
        {Tijdschrift voor Verpleegkunde}
        {\url{https://www.tvz.nl}}
        {}
        {\begin{itemize}
\item Artikel over knelpunten en verbeteringen bij medicatieoverdracht tussen ziekenhuis en thuiszorg
\item Gebaseerd op een kwaliteitsproject binnen de Sint-Maartens Kliniek
\end{itemize}}
\section{Projecten}

\cventry{jan 2022–dec 2022}
        {Implementatie van een vernieuwd medicatieprotocol om overdrachtsfouten te verminderen.}
        {Medicatieprotocol revalidatieafdeling}
        {\url{https://www.sintmaartenskliniek.nl}}
        {}
        {\begin{itemize}
\item Analyseerde incidentmeldingen en observeerde medicatieprocessen
\item Ontwikkelde samen met apotheek en artsen een checklist voor medicatieoverdracht
\item Verzorgde trainingen aan verpleegkundigen en verzorgenden
\item Leidde tot 25\% minder geregistreerde medicatie-incidenten
\end{itemize}
\textbf{Trefwoorden}: Kwaliteitsverbetering, Patiëntveiligheid, Medicatie, Implementatie}
\section{Interesses}

\cvline{Gezondheid en beweging}{Hardlopen, Yoga, Wandelen}
\cvline{Vrijwilligerswerk}{Buurtzorg, Eenzaamheidspreventie}
\cvline{Creatief}{Koken, Fotografie}
\section{Vrijwilligerswerk}

\cventry{sep 2019–jun 2024}
        {Vrijwilliger EHBO}
        {Rode Kruis Utrecht}
        {\url{https://www.rodekruis.nl}}
        {}
        {\begin{itemize}
\item Verleende EHBO-hulp bij lokale evenementen en sportwedstrijden
\item Hield EHBO-kennis en vaardigheden actueel door jaarlijkse herhalingstrainingen
\item Droeg bij aan voorlichting over eerste hulp en gezondheid in de wijk
\end{itemize}}
    
\end{document}